% Insert once after ordinary-propagation.tex.  This file supplies scalar
% comparisons only; the finite model and the signed exponential estimate
% are unchanged.  It does not use any local terminal certificate.
\paragraph{Scalar comparisons for ordinary moments}
\label{part:oc:scalar-comparisons}

\begin{lemma}[Quadratic Hermite comparison]\label{lem:oc:hermite}
Let \(W\) be a finite real variable with \(|W|\le U\), and let
\(f\in C^3([-U,U])\) satisfy \(f'''\ge0\).
For \(-U\le c<U\), put
\begin{equation}\label{eq:oc:hermite-majorant}
 Q_c(w)=f(c)+f'(c)(w-c)+b_c(w-c)^2,\qquad
 b_c=\frac{f(U)-f(c)-f'(c)(U-c)}{(U-c)^2}.
\end{equation}
Then \(f(w)\le Q_c(w)\) throughout \([-U,U]\).
If also \(f',f''\ge0\), \(0\le a\le U\),
\(\E W\le-a\) and \(\E W^2\le v_0\), then
\begin{equation}\label{eq:oc:negative-contact}
 \E f(W)\le f(-a)+b_{-a}(v_0-a^2).
\end{equation}
Here the case \(U=a=0\) is read directly. If in addition \(f\ge0\),
then \(0\le b_{-a}\le f(U)/(U+a)^2\) when \(U+a>0\).
In particular, for \(0\le h\le1\),
\begin{equation}\label{eq:oc:fixed-contact}
 \E W\le0,\quad \E W^2\le hU^2
 \quad\Longrightarrow\quad
 \E f(W)\le\frac{f(-hU)+h f(U)}{1+h},
 \qquad(f',f'',f'''\ge0).
\end{equation}

For moment matching, write \(m=\E W\), \(v=\E(W-m)^2\),
\(M=U-m>0\). With
\begin{equation}\label{eq:oc:matched-contacts}
 c=m-v/M,\qquad \ell=M+v/M,\qquad
 \pi=\frac{v}{M^2+v},
\end{equation}
one has \(-U\le c\le m<U\), and the Bernoulli variable
\(Z\) of mean \(\pi\) gives
\begin{equation}\label{eq:oc:hermite-matching}
 \E f(W)\le \E f\bigl(m+\ell(Z-\pi)\bigr).
\end{equation}
If \(M=0\), then \(W=U\); if \(v=0\), both sides equal \(f(m)\).
The matching statement needs only \(f'''\ge0\), not a restriction
on the sign of \(m\) or a bound involving \(\theta\).
\end{lemma}
\begin{proof}
The Hermite remainder has sign
\(f'''(\xi)(w-c)^2(w-U)/6\le0\).
Under the additional hypothesis \(f''\ge0\), convexity gives
\(b_c\ge0\); at \(c=-a\), the coefficient of
\(\E W\) in \(\E Q_c(W)\) is \(f'(c)+2ab_c\ge0\).
Substituting the two moment bounds proves
\eqref{eq:oc:negative-contact}; the upper bound for \(b_{-a}\)
follows by dropping the nonnegative terms in its numerator.
For \eqref{eq:oc:fixed-contact}, take \(c=-hU\): the enlarged
moments \(0,hU^2\) are realized at \(-hU,U\) with probabilities
\(1/(1+h),h/(1+h)\).
For matching, \(v\le U^2-m^2=M(U+m)\) places the lower contact
in the interval. The two-point variable in
\eqref{eq:oc:hermite-matching} has the same first two moments
as \(W\) and lies at the contact points, so its expectation of
\(Q_c\) is \(\E Q_c(W)\).
\end{proof}

The centered Bernoulli estimate used with these comparisons is
\eqref{eq:Q102:3.5}, proved in paragraph~\ref{part:Q81.2}.
Moment matching retains the actual mean of \(W\), of either sign.
The continuous balanced coefficient is \(1/8\); at a deterministic
endpoint the centered variable is zero.

\begin{lemma}[Ordinary moments from a positive sum or square]
\label{lem:oc:ordinary-comparison}
Use the original finite kernel and \(B=-K\).
For \(r\in\{4,8\}\), write \(C_r=17r^2/64\).
If \(H\ge0\), \(a>0\), and
\(H\le a^2(c+BV/2)^2\) with \(c+BV/2\ge0\), then
\begin{equation}\label{eq:oc:ordinary-square}
 \log\E H^r\le
 2r\log a+2r(c-1)+r\E BV+C_r\|V\|_2^2.
\end{equation}
If \(X,Y\ge0\), \(X+Y\le2+c+BV\), and
\(0\le H\le a^2e^{-\Gamma}XY\), then
\begin{equation}\label{eq:oc:ordinary-sum}
 \log\E H^r\le
 r(2\log a-\Gamma+c+\E BV)+C_r\|V\|_2^2.
\end{equation}
Both formulas allow either sign of \(\E BV\) and a constant
coordinate in \(V\). A zero moment is read with \(\log0=-\infty\).
\end{lemma}
\begin{proof}
For \(y\ge0\), \(y^{2r}\le e^{2r(y-1)}\).
Apply \eqref{eq:oc:signed-moments} to the resulting exponential
of \(rBV\) to obtain \eqref{eq:oc:ordinary-square}.
For the sum form, AM--GM gives
\(XY\le(1+(c+BV)/2)^2\), whose base is nonnegative since it
dominates \((X+Y)/2\); use the square form with
\(a e^{-\Gamma/2}\) and \(1+c/2\).
\end{proof}
